\documentclass[10pt,a4paper]{article}
\usepackage[T1]{fontenc}
\usepackage[utf8]{inputenc}
\usepackage{lmodern}
\usepackage[margin=2cm]{geometry}
\usepackage{booktabs}
\usepackage{tabularx}
\usepackage{xcolor}
\usepackage{enumitem}
\usepackage{hyperref}
\usepackage{array}
\hypersetup{colorlinks=true,urlcolor=blue,linkcolor=black}
\setlist{itemsep=1pt,topsep=2pt}
\newcolumntype{L}{>{\raggedright\arraybackslash}X}
\newcommand{\ok}{\textcolor{green!50!black}{\textbf{Kapandı}}}
\newcommand{\limited}{\textcolor{orange!80!black}{\textbf{Kapandı (sınırlı)}}}
\newcommand{\q}[1]{\textit{``#1''}}
\setlength{\parskip}{4pt}
\setlength{\parindent}{0pt}

\title{\textbf{CNSM 2026 Hakem Eleştirileri ve Çözümleri}\\[2pt]
\large DEMO: From Natural Language to Live 5G API Calls\\
\normalsize CNSM 2026 demo \#1571315592 $\rightarrow$ WMNC 2026 demo yeniden gönderimi}
\author{Türkan Doğa Durak}
\date{13 Ekim 2026}

\begin{document}
\maketitle

\section*{Özet}
CNSM 2026'da dört hakemin dördü de \textit{Weak Accept (3)} verdi; konu uygunluğu tüm hakemlerce
\q{Very relevant (4)} olarak puanlandı. Ret kararı kavramsal değil, büyük ölçüde \textbf{doğrulanabilirlik ve
tutarlılık} sorunlarından kaynaklandı: makaledeki ``canlı'' iddiaları ile şekil/kod arasındaki çelişkiler,
public repodaki eski sonuçlar ve değerlendirmenin kapsamının belirsizliği. Toplam 13 eleştirinin 10'u tamamen,
3'ü sınırlamaları açıkça belirtilerek kapatıldı. Çözümlerin çoğu yalnızca metin düzeltmesi değil; kodda,
test ortamında ve verilerde gerçek değişiklik içeriyor. Tüm kanıtlar
\url{https://github.com/doadurak/mini-telco-platform} adresinde.

\begin{table}[h]
\centering\small
\begin{tabularx}{\textwidth}{l L l}
\toprule
\textbf{Hakem} & \textbf{Ana kaygı} & \textbf{Durum} \\
\midrule
H1 & ``Canlı'' denilen Fig.~2'de \texttt{dry\_run: true}; değerlendirmenin neyi ölçtüğü belirsiz & 4/4 \ok \\
H2 & İlgili çalışmalarla deneysel karşılaştırma, ölçeklenebilirlik, LLM bağımlılığı, biçim & 6/6 (2'si sınırlı) \\
H3 & Katman bazlı zorluklar; doğruluk ve gecikmeye etkisi & 1/1 \ok \\
H4 & Public repo makaleyle çelişiyor (60-entry, Gemini 6/60, tarihsiz bug fix) & 3/3 \ok \\
\bottomrule
\end{tabularx}
\end{table}

\section{Hakem 1: Dry-run ile canlı yürütmenin karışması}
\textbf{Neye itiraz etti?} Makale Fig.~2'yi ``canlı provider yanıtı ve üç doğrulama katmanı'' olarak anlatıyordu,
ama şekildeki istekte \texttt{"dry\_run": true} vardı: \q{The authors should clarify whether this execution includes a
provider call and a live CAPIF registry check.} Hakem ayrıca yanlış ama kayıtlı bir servisin nasıl yakalandığını ve
300 kaydın ne olduğunu sordu.

\renewcommand{\arraystretch}{1.15}
\begin{tabularx}{\textwidth}{>{\raggedright\arraybackslash}p{3.6cm} L l}
\toprule
\textbf{Eleştiri} & \textbf{Ne yaptık / kanıt} & \textbf{Durum} \\
\midrule
Fig.~2 \texttt{dry\_run: true} gösteriyor &
Fig.~2, gerçek 5G core üzerinde yapılmış bir \texttt{dry\_run:false} çağrısının kayıtlı yanıtından yeniden üretildi:
QoD HTTP~201, Open5GS PCF onayıyla \texttt{qosStatus: AVAILABLE}, \texttt{registry\_checked=true}.
Kanıt: \texttt{docs/live\_runs/2026-10-11-5gc/}. Repodaki eski görselin aslında CAPIF bağlantısı \emph{başarısızken}
(\texttt{mock-onboarded}) alındığı da tespit edilip değiştirildi. & \ok \\
Provider çağrısı ve canlı CAPIF kontrolü var mı? &
\textbf{Kodda düzeltildi:} L3 eskiden \texttt{dry\_run} açıkken atlanıyordu (\texttt{skip\_registry=dry\_run}).
Artık L3 her zaman canlı CAPIF registry'sine soruyor (\texttt{check\_registry}). mTLS'in yalnızca CAPIF
onboarding/discovery'de, provider çağrısının HTTP ile yapıldığı açıkça yazıldı. & \ok \\
Yanlış ama kayıtlı servis nasıl yakalanıyor? &
Dürüst cevap: yakalanmıyor. L3 yalnızca servisin yayınlı olduğunu kontrol eder; feedback döngüsü yalnızca
payload'ı düzeltir. Hata enjeksiyonu bunu sayıyla gösteriyor: L3 hataları 0/20 düzeltildi. S3 senaryosu bu
sınırı gösteren bir demo örneğine dönüştürüldü. & \ok \\
300 kayıt, ground truth, IWSR'nin kapsamı &
300 \emph{benzersiz} intent (152 EN / 148 TR), intent başına tek koşu, ground truth CAMARA OpenAPI kısıtları ve
NEF test verisinden. Benchmark provider çağırmaz; kayıtlı 1.200 planın tamamı canlı L3'e karşı yeniden doğrulandı
(0 ret). & \ok \\
\bottomrule
\end{tabularx}

\section{Hakem 2: Karşılaştırma, ölçeklenebilirlik, sunum}
\textbf{Neye itiraz etti?} \q{The comparison with related work is mostly qualitative and lacks experimental
benchmarking}; demo bölümünün çok sayıda alt bölüme bölünmesi; gecikme kırılımı, ölçeklenebilirlik ve harici LLM
API'lerine bağımlılığın yüzeysel kalması; yüzde boşlukları ve bölünen kelimeler.

\begin{tabularx}{\textwidth}{>{\raggedright\arraybackslash}p{3.6cm} L l}
\toprule
\textbf{Eleştiri} & \textbf{Ne yaptık / kanıt} & \textbf{Durum} \\
\midrule
Deneysel karşılaştırma yok &
İki ablasyon (GPT-4o-mini, 300 kayıt): \textbf{şemasız LLM IWSR \%16.3} (model alan adlarını uyduruyor;
Liu ve ark.'nın grounding'siz GPT-4 için \%40'ı ile karşılaştırıldı), \textbf{tek çağrılı LLM \%77.3},
bizim iki fazlı planlayıcımız \%96.7. Tablo~I'e eklendi. Harici bir sistemle doğrudan benchmark yapılmadı. & \limited \\
Çok fazla alt bölüm & Senaryolar tek bölümde, kalın başlıklı paragraflar halinde birleştirildi. & \ok \\
Gecikme kırılımı & Her yanıtta aşama bazlı süre: planlama toplamın \%78--99'u, doğrulama 13--27~ms. & \ok \\
Ölçeklenebilirlik &
Yük testi: LLM'siz platform $\approx$160~istek/s (20 istemcide p95 156~ms); GPT-4o-mini ile gecikme
$\approx$2.6~s'de sabit, verim 0.34$\rightarrow$3.3~istek/s doğrusal. Darboğaz LLM API'si. Tek host üzerinde
ölçüldü; dağıtık ölçekleme denenmedi. & \limited \\
Harici LLM API bağımlılığı &
API hatasında \texttt{auto} modu kural tabanlı planlayıcıya dönüyor (bu davranıştaki bir hata da düzeltildi:
yalnızca Gemini anahtarı varsa çalışıyordu). MSISDN'lerin dış API'ye gittiği ve üretimde pseudonymization veya
self-hosted model gerektiği yazıldı. & \ok \\
Biçim tutarsızlıkları &
\% öncesi boşluk kaldırıldı, CAMARA/OpenCAPIF artık satır sonunda bölünmüyor, aralıklarda en-dash. Derleme logu
uyarısız. Danışmanın ``not clear sentence'' yorumu üzerine sıkıştırılmış cümleler açık hale getirildi. & \ok \\
\bottomrule
\end{tabularx}

\section{Hakem 3: Katman bazlı zorluklar}
\textbf{Neye itiraz etti?} \q{This demo could be expanded by preparing a few test sets to illustrate the
difficulties encountered at the various layers \ldots\ and by presenting the impact on the accuracy and also on
other metrics e.g., latency.}

\begin{tabularx}{\textwidth}{>{\raggedright\arraybackslash}p{3.6cm} L l}
\toprule
\textbf{Eleştiri} & \textbf{Ne yaptık / kanıt} & \textbf{Durum} \\
\midrule
Katman bazlı test setleri, doğruluk ve gecikmeye etkisi &
Doğru planlara 140 kontrollü hata enjekte edildi (Tablo~II): hepsi hedef katmanda yakalandı. Feedback L1'de
40/40, L2'de 79/80 hatayı tek turda ($\approx$1.3~s) düzeltti; L3'te 0/20 (3 tur, $\approx$3.8~s). Ayrıca 300
kayıtta hata türü (yanlış servis, L1/L2 reddi) ve zorluk bazlı IWSR (zor: \%87--95, kolay/orta: \%99--100)
raporlandı. L3 hatası registry yanıtından servis çıkarılarak simüle edildi. & \ok \\
\bottomrule
\end{tabularx}

\section{Hakem 4: Public repo ile makale arasındaki çelişki}
\textbf{Neye itiraz etti?} En sert ve en somut eleştiri buydu. Hakem repoyu incelemiş ve
\q{this is directly contradicted by the authors' own public repository} demişti: README'de 60 kayıtlık eski
tabloda deterministik baseline \%100 ile LLM'lerin önündeydi, Gemini \q{partial 6/60-entry run} görünüyordu ve
bug fix'ler veri toplama tarihine göre tarihsizdi.

\begin{tabularx}{\textwidth}{>{\raggedright\arraybackslash}p{3.6cm} L l}
\toprule
\textbf{Eleştiri} & \textbf{Ne yaptık / kanıt} & \textbf{Durum} \\
\midrule
60-entry ters sıralama &
README ve demo script'i 300 kayıtlık sonuçları gösteriyor; pilot \texttt{datasets/legacy\_pilot\_60/}
altında arşivlendi. Çelişkinin \emph{sebebi} ölçüldü: anahtar kelime kuralları pilot için yazıldığından baseline
pilotta 60/60, görülmemiş 240 kayıtta \%79.2; LLM'ler aynı 240 kayıtta \%95.4--97.1. & \ok \\
Gemini 6/60 &
Gemini'nin pilotu 60/60 tamamlanmıştı, 300/300 sonuçları da var; README ve demo script'indeki eski ifadeler
temizlendi. & \ok \\
Tarihsiz bug fix'ler &
README'de git geçmişine dayalı tarihli zaman çizelgesi; pilot sonuçları üç hatanın izleri için tarandı, etkilenen
kayıt yok. Çok-niyetli cümlelerdeki bir planlayıcı hatası da bu süreçte bulunup düzeltildi ve etkilenen 11 kayıt
izlenebilir şekilde yeniden koşturuldu (sonuç dosyalarında \texttt{rerun} alanı). & \ok \\
\bottomrule
\end{tabularx}

\section{Eleştirilerin ötesinde yapılan iyileştirmeler}
\begin{itemize}
\item \textbf{Gerçek 5G core:} Open5GS 5GC + UERANSIM (3 UE) eklendi. QoS istekleri artık
CAMARA $\rightarrow$ NEF $\rightarrow$ Open5GS PCF zincirinde gerçekten yetkilendiriliyor (\texttt{testbed/open5gs}).
\item \textbf{Profil uyumsuzluğu:} LLM'in ürettiği \texttt{QOS\_E\_STREAMING} gibi profilleri provider
tanımıyordu; operatör profillerine eşleme eklendi.
\item \textbf{Telefon numarası $\rightarrow$ UE IP çözümlemesi} 5GC abone veritabanından yapılıyor.
\item \textbf{Yeni mimari şekli (Fig.~1)} kodla birebir doğrulandı (CAPIF-1e/mTLS, CAPIF-2e/JWT, CAPIF-4).
\item \textbf{Yeniden üretilebilirlik:} \texttt{start.sh}, \texttt{scripts/reset\_capif.sh},
\texttt{scripts/run\_live\_scenarios.sh}; test paketi 78/78.
\end{itemize}

\section{Açıkça belirtilen sınırlamalar}
\begin{itemize}
\item UERANSIM NGAP \textit{PDU Session Resource Modify} desteklemediği için QoS akışı core tarafından
yetkilendiriliyor ama emüle radyoda uygulanmıyor.
\item Konum verisi NEF'in emüle veri kaynağından geliyor, Open5GS AMF'den değil.
\item Ablasyonlar yalnızca GPT-4o-mini ile, intent başına tek koşu yapıldı.
\item Ölçeklenebilirlik tek host üzerinde ölçüldü; L3 hata enjeksiyonu simülasyon.
\end{itemize}

\end{document}
